\documentclass[../main.tex]{subfiles}
\begin{document}

\section{积性函数前缀和}

\subsection{杜教筛}

\subsubsection{块筛}

\begin{frame}
记 $S_f(n)=\sum\limits_{i=1}^nf(i)$。块筛定义为：求出
$$
\{(\lfloor\frac{n}{i}\rfloor,S_f(\lfloor\frac{n}{i}\rfloor))\mid i\in\{1,2,\dots,n\} \}
$$
集合大小是 $O(\sqrt n)$ 的。
\end{frame}

\subsubsection{杜教筛}

\begin{frame}
杜教筛的核心思想是寻找 $f=g/h$，满足 $g,h$ 可块筛，则可以用 $O(n^{2/3})$ 复杂度块筛 $f$。
$$
\begin{aligned}
S_g(n)&=\sum_{i=1}^n\sum_{d\mid i}h(d)f(\frac{i}{d})\\
&=\sum_{d=1}^nh(d)S_f(\lfloor\frac{n}{d}\rfloor)\\
&=S_f(n)h(1)+\sum_{d=2}^nh(d)S_f(\lfloor\frac{n}{d}\rfloor)\\
S_f(n)&=\frac{1}{h(1)}(S_g(n)-\sum_{d=2}^nh(d)S_f(\lfloor\frac{n}{d} \rfloor))
\end{aligned}
$$
\end{frame}

\begin{frame}
上式可以整除分块，枚举每个 $m=\lfloor\dfrac{n}{d}\rfloor$，满足 $\lfloor\dfrac{n}{d}\rfloor=m$ 的 $d$ 为 $\lfloor\dfrac{n}{m+1}\rfloor+1\sim\lfloor\dfrac{n}{m}\rfloor$。由于 $h$ 可块筛，故 $S_f(n)$ 可以 $O(\sqrt n)$ 计算。

分析复杂度。在每个 $n$ 处复杂度都是 $O(\sqrt n)$，总复杂度为 $O(\sum\limits_{i=1}^{\sqrt n}\sqrt i+\sqrt{n/i})=O(n^{3/4})$。

如果使用线筛预处理出前 $n^{2/3}$ 项，则复杂度降为 $O(\sum\limits_{i=1}^{n^{1/3}}\sqrt{n/i})=O(n^{2/3})$。

实际计算中可以在计算 $S_f(n)$ 时递归计算每个 $S_f(\lfloor\dfrac{n}{d}\rfloor)$。求出 $S_f(n)$ 的同时得出所有块筛。
\end{frame}

\subsubsection{积式}

\begin{frame}
如果不是 $f=g/h$ 而是 $f=g*h$，同样可以求 $S_f(n)$ 的块筛。
$$
\begin{aligned}
S_f(n)&=\sum_{i=1}^n\sum_{d\mid i}g(d)h(\frac{i}{d})\\
&=\sum_{d=1}^ng(d)S_h(\lfloor\frac{n}{d}\rfloor)
\end{aligned}
$$
如果已知 $g,h$ 的块筛，则对 $S_f$ 求单点需要 $O(\sqrt n)$，求块筛同样可以预处理前 $n^{2/3}$ 项达到 $O(n^{2/3})$。
\end{frame}

\subsubsection{练习}

\begin{frame}
用 $O(n^{2/3})$ 复杂度求以下函数的块筛：

\begin{itemize}
\item $\varphi$
\item $\mu$
\item $\sigma_k$
\item $\varphi\cdot \text{id}_k$
\item $d\cdot \text{id}_k$
\end{itemize}

提示：使用贝尔级数来找合适的 $f=g/h$。
\end{frame}

\subsubsection{P4318 完全平方数}

\begin{frame}
$\mu^2(n)=\mu(n)\cdot\mu(n)$。求 $S_{\mu^2}(n)$。

$n$ 可以被唯一地分解为 $x^2y$，其中 $\mu^2(y)=1$。

令 $f(n)=[\exists d,d^2=n]$，则 $1=f*\mu^2$。

$S_f(n)=\lfloor \sqrt n\rfloor$ 可以 $O(1)$ 计算，故可以杜教筛。复杂度 $O(n^{2/3})$。
\end{frame}

\begin{frame}
更优秀的做法：
$$
\begin{aligned}
\mu^2(n)=[x=1]=\sum_{d\mid x}\mu(d)=\sum_{d^2\mid n}\mu(d)\\
S_{\mu^2}(n)=\sum_{i=1}^n\sum_{d^2\mid i}\mu(d)=\sum_{d=1}^{\sqrt{n}}\mu(d)\lfloor\frac{n}{d^2}\rfloor
\end{aligned}
$$
复杂度 $O(\sqrt n)$。
\end{frame}

\subsubsection{来源不明题}

\begin{frame}
将 $n$ 质因数分解为 $\sum\limits_{i=1}^kp_i^{\alpha_i}$，定义 $f(n)=(-1)^{\sum_{i=1}^k\alpha_i}$。杜教筛 $S_f$。

显然 $f$ 是积性函数。

计算其贝尔级数 $f_p(x)=\sum\limits_{i=0}^{\infty}(-1)^ix^i=\dfrac{1}{1+x}$。于是 $f=\epsilon/\mu^2$。由上题，$\mu^2$ 可块筛。故 $S_f$ 可以 $O(n^{2/3})$ 求。
\end{frame}

\subsubsection{51nod1220 约数之和}

\begin{frame}
求 $\sum\limits_{i=1}^n\sum\limits_{j=1}^n\sigma(ij)$。$n\le 10^9$。
\end{frame}

\begin{frame}
\begin{theo}[引理]
$$
\sigma(xy)=\sum_{i\mid x}\sum_{j\mid y}\frac{x}{i}\cdot j[\gcd(i,j)=1]
$$
\end{theo}

\begin{proof}
在集合 $S=\{(i,j)\mid i|x,j| y,\gcd(i,j)=1 \}$ 与 $D=\{d\mid xy\}$ 之间构造映射。$(i,j)$ 对应于 $d=\dfrac{x}{i}\cdot j$，$d$ 对应于 $(\dfrac{x}{\gcd(d,x)},\dfrac{d}{\gcd(d,x)})$。

分别在每个素因子上考虑，可以发现两个映射都是单射，且互为逆映射。故为双射。
\end{proof}
\end{frame}

\begin{frame}
根据引理，原式化为
$$
\begin{aligned}
&\sum_{x=1}^n\sum_{y=1}^n\sum_{i\mid x}\sum_{j\mid y}\frac{x}{i}j\sum_{d\mid i,j}\mu(d)\\
=&\sum_{d=1}^n\mu(d)(\sum_{d\mid i}\sum_{i\mid x}\frac{x}{i})(\sum_{d\mid j}\sum_{j\mid y}j)\\
=&\sum_{d=1}^n\mu(d)(\sum_{d\mid x}\sum_{d\mid i\mid x}\frac{x}{i})(\sum_{d\mid y}\sum_{d\mid j\mid y}d\cdot\frac{j}{d})\\
=&\sum_{d=1}^n\mu(d)\cdot d\cdot S_{\sigma}^2(\lfloor\frac{n}{d}\rfloor)
\end{aligned}
$$
求 $\mu\cdot \text{id}$ 和 $\sigma$ 的块筛即可。复杂度 $O(n^{2/3})$。
\end{frame}

\subsubsection{P1587 [NOI2016] 循环之美}

\begin{frame}
给定 $n,m,k$，求在 $k$ 进制下，有多少个数值上互不相等的纯循环小数，可以用分数 $\dfrac{x}{y}$ 表示，其中 $1\le x\le n,1\le y\le m$，且 $x,y\in\Z$。一个数是纯循环的，当且仅当它在 $k$ 进制下可以写成 $a.\dot{c_1}c_2c_3\dots c_{l-1}\dot{c_l}$。

$n,m\le 10^9,k\le 2000$。
\end{frame}

\begin{frame}
设 $\dfrac{x}{y}$ 是纯循环小数，循环节长度为 $l$。有 $\dfrac{xk^l}{y}-\dfrac{x}{y}\in\Z$，即 $xk^l\equiv x\pmod y$。而 $\gcd(x,y)=1$，故 $k^l\equiv 1\pmod y$，等价于 $\gcd(k,y)=1$。易知这也是充分条件。

故答案为
$$
\sum\limits_{x=1}^n\sum\limits_{y=1}^m[\gcd(x,y)=1][\gcd(y,k)=1]
$$
\end{frame}

\begin{frame}
$$
\begin{aligned}
\text{ans}=&\sum_{d=1}^n\mu(d)\sum_{d\mid x}\sum_{d\mid y}[\gcd(y,k)=1]\\
=&\sum_{d=1}^n\mu(d)\cdot\chi_k(d)\cdot\lfloor\frac{n}{d}\rfloor\cdot S_{\chi_k}(\lfloor\frac{m}{d}\rfloor)
\end{aligned}
$$
其中 $\chi_k(x)=[\gcd(x,k)=1]$，是一个完全积性函数。$S_{\chi_k}(x)=\lfloor\dfrac{x}{k}\rfloor\varphi(k)+S_{\chi_k}(x\bmod k)$，可以 $O(k)$ 预处理后 $O(1)$ 查询。

于是只需再块筛 $\mu\cdot\chi_k$。$(\mu\cdot\chi_k)*\chi_k=(\mu\cdot\chi_k)*(1\cdot\chi_k)=(\mu*1)\cdot\chi_k=\epsilon$。杜教筛即可。复杂度 $O(n^{2/3})$。
\end{frame}

\subsection{PN 筛}

\subsubsection{Powerful Number}

\begin{frame}
将 $n$ 质因数分解，若每个素数的指数均大于等于 $2$，则称 $n$ 是一个 Powerful Number。

\begin{theo}[结论]
$n$ 以内的 PN 数量为 $O(\sqrt n)$ 级别。
\end{theo}

\begin{proof}
一个 PN 可以被分解为 $a^2b^3$ 的形式。
$$
O(\sum_{a=1}^{\sqrt n}\sqrt[3]{\frac{n}{a^2}})=O(n^{1/3}\int_{1}^{n^{1/2}}x^{-2/3}\mathrm{d}x)=O(\sqrt n)
$$
\end{proof}
\end{frame}

\begin{frame}
如何找 $n$ 以内所有 PN？取出 $\sqrt n$ 以内所有素数，用 DFS 对这些素数及其素数幂进行组合即可。复杂度 $O(\sqrt n)$。
\end{frame}

\subsubsection{素数拟合}

\begin{frame}
给定一个积性函数 $f$，求 $S_f(n)$。

做法：构造一个易块筛的积性函数 $g$，满足 $g$ 在所有素数处都与 $f$ 相等。称为 $g$ 与 $f$ 素数拟合。

再令 $h=f/g$，此时 $h$ 也是积性函数。由 $f(p)=g(p)h(1)+g(1)h(p),g(p)=f(p)$，知 $h(p)=0$。于是 $h$ 只在 PN 处有值。
$$
\begin{aligned}
S_f(n)&=\sum_{i=1}^n\sum_{d\mid i}h(d)g(\frac{i}{d})\\
&=\sum_{d\in \text{PN}}h(d)S_g(\lfloor\frac{n}{d}\rfloor)
\end{aligned}
$$
\end{frame}

\begin{frame}
接下来求 $h$。因为 $h$ 有积性，故只需要求出其在素数幂处的取值。
$$
\begin{aligned}
f(p^k)=\sum_{i=0}^kh(p^i)g(p^{k-i})\\
h(p^k)=f(p^k)-\sum_{i=1}^{k-1}h(p^i)g(p^{k-i})
\end{aligned}
$$
所有的 $h(p^k)$ 都可以在找 PN 的过程中暴力求出。这部分的复杂度为 $O(\sum\limits_{p\in\mathbb{P},p\le \sqrt n}\log_p^2(n))=O(\dfrac{\sqrt n}{\log n})$，不为瓶颈。
\end{frame}

\begin{frame}
一般复杂度瓶颈在于块筛 $g$。常为 $O(n^{2/3})$。

特别地，如果 $S_g(x)$ 能在 $O(\sqrt x)$ 内求出，则复杂度优化至
$$
O(\sum_{d\in\text{PN}}\sqrt{\frac{n}{d}})=O(\sum_{a,b}\sqrt{\frac{n}{a^2b^3}})=O(n^{1/2}\sum_a\frac{1}{a})=O(\sqrt n\log n)
$$
\end{frame}

\subsubsection{P5325 【模板】Min\_25 筛}

\begin{frame}
定义积性函数 $f$，$f(p^k)=p^k(p^k-1)$。求 $S_f(n)$。$n\le 10^{10}$。

令 $g=\text{id}\cdot\varphi$。则 $g$ 与 $f$ 素数拟合。

$\text{id}\cdot\varphi=\text{id}\cdot(\text{id}/1)=\text{id}_2/\text{id}$，可以使用杜教筛求出块筛。复杂度 $O(n^{2/3})$。
\end{frame}

\subsubsection{LOJ6053 简单的函数}

\begin{frame}
定义积性函数 $f$，$f(p^c)=p\oplus c$，其中 $\oplus$ 为异或。求 $S_f(n)$。$n\le10^{10}$。
\end{frame}

\begin{frame}
当 $p$ 为 $2$ 时，$f(p)=3$。当 $p$ 为奇素数时，$f(p)=p-1$。

于是构造 $g(x)=\left\{\begin{aligned}&3\varphi(x) & 2\mid x\\ &\varphi(x) & 2\nmid x\end{aligned}\right.$，则 $g$ 为积性函数且与 $f$ 素数拟合。

只需块筛 $S_g(n)=\sum\limits_{i=1}^n\varphi(i)+2\sum\limits_{i=1}^{\lfloor n/2\rfloor}\varphi(2i)$。第一部分 $S_{\varphi}$ 可以用杜教筛。

第二部分 $F(n)=\sum\limits_{i=1}^n\varphi(2i)=\sum\limits_{i\le m,2\mid i}2\varphi(i)+\sum\limits_{i\le m,2\nmid i}\varphi(i)=S_{\varphi}(n)+\sum\limits_{i=1}^{\lfloor n/2\rfloor}\varphi(2i)$。于是得到递推式 $F(n)=S_{\varphi}(n)+F(\lfloor n/2\rfloor)$。递归计算的复杂度为 $O(\log n)$，不为瓶颈。

总复杂度为块筛 $S_{\varphi}$ 的 $O(n^{2/3})$。
\end{frame}

\end{document}
